\subsection{Análisis de tráfico de los seis métodos trace}\label{sec:trace-patrones}
Las siguientes seis fichas usan Valparaíso como ejemplo homogéneo y conservan comprobaciones de los tres destinos en la evidencia digital. Los valores en hexadecimal son bytes de aplicación; los encabezados TCP/UDP/ICMP se analizan por separado.

\Needspace{9\baselineskip}
\subsubsection{scamper / trace / udp}
Ejemplo: Valparaíso. Entrenamiento: scamper\_\allowbreak{}20261026T204619\_\allowbreak{}270c47c3; contraste: rondas 2 y 3. Las reglas completas de los tres destinos están en el archivo de predicciones enlazado en el anexo.

\begingroup
\fontsize{9}{11}\selectfont
\addtocounter{table}{-1}
\begin{longtable}{p{3.1cm}p{5.8cm}p{5.1cm}}
\toprule
\textbf{Campo} & \textbf{Regla en ronda 1} & \textbf{Contraste posterior} \\ \midrule
\endfirsthead
\toprule
\textbf{Campo} & \textbf{Regla en ronda 1} & \textbf{Contraste posterior} \\ \midrule
\endhead
ip.ttl & incremento 1 por paquete & r2: cumple; r3: cumple \\
ip.id & incremento 1 por paquete & r2: cumple; r3: cumple \\
ip.len & constante 44 & r2: cumple; r3: cumple \\
ip.proto & constante 17 & r2: cumple; r3: cumple \\
udp.srcport & constante 32769 & r2: cumple; r3: cumple \\
udp.dstport & incremento 1 por paquete & r2: cumple; r3: cumple \\
udp.length & constante 24 & r2: cumple; r3: cumple \\
udp.checksum & incremento -1 por paquete & r2: cumple; r3: cumple \\
\bottomrule
\end{longtable}
\endgroup
Payload: [16] bytes. Hexadecimal inicial (offset 0):

\begin{quote}\ttfamily\footnotesize\raggedright 0000 0000 0000 0000 0000 0000 0000 0000\end{quote}
La suma de payload\_\allowbreak{}word0\_\allowbreak{}BE, udp.checksum, udp.srcport, udp.dstport módulo 65535 se mantuvo constante. La base se aprende del primer paquete de cada ejecución; no se presume universal. Es compatible con compensación de checksum, sin probar la intención del emisor.

Predicción concreta en ronda 2, campo udp.dstport: la primera solicitud tiene valor 33435. Con la regla aprendida y esa base cuando corresponde, se predicen las solicitudes 2 a 4 antes de contrastarlas con sus valores observados.

\begingroup
\fontsize{9}{11}\selectfont
\addtocounter{table}{-1}
\begin{longtable}{p{3.2cm}p{5.1cm}p{5.1cm}}
\toprule
\textbf{Solicitud / trama} & \textbf{Predicción} & \textbf{Observación} \\ \midrule
\endfirsthead
\toprule
\textbf{Solicitud / trama} & \textbf{Predicción} & \textbf{Observación} \\ \midrule
\endhead
2 / 7 & 33436 & 33436 \\
3 / 9 & 33437 & 33437 \\
4 / 11 & 33438 & 33438 \\
\bottomrule
\end{longtable}
\endgroup
Alcance: 9 solicitudes en el ejemplo inicial. Una verificación de campo o byte no equivale a una medición independiente; los campos variables conservan sus restricciones y no reciben valores inventados.

\Needspace{9\baselineskip}
\subsubsection{scamper / trace / icmp}
Ejemplo: Valparaíso. Entrenamiento: scamper\_\allowbreak{}20261026T204623\_\allowbreak{}0cf6d8d5; contraste: rondas 2 y 3. Las reglas completas de los tres destinos están en el archivo de predicciones enlazado en el anexo.

\begingroup
\fontsize{9}{11}\selectfont
\addtocounter{table}{-1}
\begin{longtable}{p{3.1cm}p{5.8cm}p{5.1cm}}
\toprule
\textbf{Campo} & \textbf{Regla en ronda 1} & \textbf{Contraste posterior} \\ \midrule
\endfirsthead
\toprule
\textbf{Campo} & \textbf{Regla en ronda 1} & \textbf{Contraste posterior} \\ \midrule
\endhead
ip.ttl & incremento 1 por paquete & r2: cumple; r3: cumple \\
ip.id & incremento 1 por paquete & r2: cumple; r3: cumple \\
ip.len & constante 44 & r2: cumple; r3: cumple \\
ip.proto & constante 1 & r2: cumple; r3: cumple \\
icmp.type & constante 8 & r2: cumple; r3: cumple \\
icmp.code & constante 0 & r2: cumple; r3: cumple \\
icmp.ident & constante 32769 & r2: cumple; r3: cumple \\
icmp.seq & incremento 1 por paquete & r2: cumple; r3: cumple \\
icmp.checksum & incremento -1 por paquete & r2: cumple; r3: cumple \\
\bottomrule
\end{longtable}
\endgroup
Payload: [16] bytes. Hexadecimal inicial (offset 0):

\begin{quote}\ttfamily\footnotesize\raggedright 0000 0000 0000 0000 0000 0000 0000 0000\end{quote}
La suma de payload\_\allowbreak{}word0\_\allowbreak{}BE, icmp.checksum, icmp.ident, icmp.seq módulo 65535 se mantuvo constante. La base se aprende del primer paquete de cada ejecución; no se presume universal. Es compatible con compensación de checksum, sin probar la intención del emisor.

Predicción concreta en ronda 2, campo icmp.seq: la primera solicitud tiene valor 0. Con la regla aprendida y esa base cuando corresponde, se predicen las solicitudes 2 a 4 antes de contrastarlas con sus valores observados.

\begingroup
\fontsize{9}{11}\selectfont
\addtocounter{table}{-1}
\begin{longtable}{p{3.2cm}p{5.1cm}p{5.1cm}}
\toprule
\textbf{Solicitud / trama} & \textbf{Predicción} & \textbf{Observación} \\ \midrule
\endfirsthead
\toprule
\textbf{Solicitud / trama} & \textbf{Predicción} & \textbf{Observación} \\ \midrule
\endhead
2 / 5 & 1 & 1 \\
3 / 7 & 2 & 2 \\
4 / 9 & 3 & 3 \\
\bottomrule
\end{longtable}
\endgroup
Alcance: 9 solicitudes en el ejemplo inicial. Una verificación de campo o byte no equivale a una medición independiente; los campos variables conservan sus restricciones y no reciben valores inventados.

\Needspace{9\baselineskip}
\subsubsection{scamper / trace / udp-paris}
Ejemplo: Valparaíso. Entrenamiento: scamper\_\allowbreak{}20261026T204627\_\allowbreak{}61f04669; contraste: rondas 2 y 3. Las reglas completas de los tres destinos están en el archivo de predicciones enlazado en el anexo.

\begingroup
\fontsize{9}{11}\selectfont
\addtocounter{table}{-1}
\begin{longtable}{p{3.1cm}p{5.8cm}p{5.1cm}}
\toprule
\textbf{Campo} & \textbf{Regla en ronda 1} & \textbf{Contraste posterior} \\ \midrule
\endfirsthead
\toprule
\textbf{Campo} & \textbf{Regla en ronda 1} & \textbf{Contraste posterior} \\ \midrule
\endhead
ip.ttl & incremento 1 por paquete & r2: cumple; r3: cumple \\
ip.id & incremento 1 por paquete & r2: cumple; r3: cumple \\
ip.len & constante 44 & r2: cumple; r3: cumple \\
ip.proto & constante 17 & r2: cumple; r3: cumple \\
udp.srcport & constante 32769 & r2: cumple; r3: cumple \\
udp.dstport & constante 33435 & r2: cumple; r3: cumple \\
udp.length & constante 24 & r2: cumple; r3: cumple \\
udp.checksum & incremento 1 por paquete & r2: cumple; r3: cumple \\
\bottomrule
\end{longtable}
\endgroup
Payload: [16] bytes. Hexadecimal inicial (offset 0):

\begin{quote}\ttfamily\footnotesize\raggedright 66b1 0000 0000 0000 0000 0000 0000 0000\end{quote}
La suma de payload\_\allowbreak{}word0\_\allowbreak{}BE, udp.checksum, udp.srcport, udp.dstport módulo 65535 se mantuvo constante. La base se aprende del primer paquete de cada ejecución; no se presume universal. Es compatible con compensación de checksum, sin probar la intención del emisor.

Predicción concreta en ronda 2, campo ip.ttl: la primera solicitud tiene valor 1. Con la regla aprendida y esa base cuando corresponde, se predicen las solicitudes 2 a 4 antes de contrastarlas con sus valores observados.

\begingroup
\fontsize{9}{11}\selectfont
\addtocounter{table}{-1}
\begin{longtable}{p{3.2cm}p{5.1cm}p{5.1cm}}
\toprule
\textbf{Solicitud / trama} & \textbf{Predicción} & \textbf{Observación} \\ \midrule
\endfirsthead
\toprule
\textbf{Solicitud / trama} & \textbf{Predicción} & \textbf{Observación} \\ \midrule
\endhead
2 / 10 & 2 & 2 \\
3 / 13 & 3 & 3 \\
4 / 15 & 4 & 4 \\
\bottomrule
\end{longtable}
\endgroup
Alcance: 9 solicitudes en el ejemplo inicial. Una verificación de campo o byte no equivale a una medición independiente; los campos variables conservan sus restricciones y no reciben valores inventados.

\Needspace{9\baselineskip}
\subsubsection{scamper / trace / icmp-paris}
Ejemplo: Valparaíso. Entrenamiento: scamper\_\allowbreak{}20261026T204633\_\allowbreak{}256949f1; contraste: rondas 2 y 3. Las reglas completas de los tres destinos están en el archivo de predicciones enlazado en el anexo.

\begingroup
\fontsize{9}{11}\selectfont
\addtocounter{table}{-1}
\begin{longtable}{p{3.1cm}p{5.8cm}p{5.1cm}}
\toprule
\textbf{Campo} & \textbf{Regla en ronda 1} & \textbf{Contraste posterior} \\ \midrule
\endfirsthead
\toprule
\textbf{Campo} & \textbf{Regla en ronda 1} & \textbf{Contraste posterior} \\ \midrule
\endhead
ip.ttl & incremento 1 por paquete & r2: cumple; r3: cumple \\
ip.id & incremento 1 por paquete & r2: cumple; r3: cumple \\
ip.len & constante 44 & r2: cumple; r3: cumple \\
ip.proto & constante 1 & r2: cumple; r3: cumple \\
icmp.type & constante 8 & r2: cumple; r3: cumple \\
icmp.code & constante 0 & r2: cumple; r3: cumple \\
icmp.ident & constante 32769 & r2: cumple; r3: cumple \\
icmp.seq & incremento 1 por paquete & r2: cumple; r3: cumple \\
icmp.checksum & constante 0x8001 & r2: cumple; r3: cumple \\
\bottomrule
\end{longtable}
\endgroup
Payload: [16] bytes. Hexadecimal inicial (offset 0):

\begin{quote}\ttfamily\footnotesize\raggedright f7fc 0000 0000 0000 0000 0000 0000 0000\end{quote}
La suma de payload\_\allowbreak{}word0\_\allowbreak{}BE, icmp.checksum, icmp.ident, icmp.seq módulo 65535 se mantuvo constante. La base se aprende del primer paquete de cada ejecución; no se presume universal. Es compatible con compensación de checksum, sin probar la intención del emisor.

Predicción concreta en ronda 2, campo icmp.seq: la primera solicitud tiene valor 0. Con la regla aprendida y esa base cuando corresponde, se predicen las solicitudes 2 a 4 antes de contrastarlas con sus valores observados.

\begingroup
\fontsize{9}{11}\selectfont
\addtocounter{table}{-1}
\begin{longtable}{p{3.2cm}p{5.1cm}p{5.1cm}}
\toprule
\textbf{Solicitud / trama} & \textbf{Predicción} & \textbf{Observación} \\ \midrule
\endfirsthead
\toprule
\textbf{Solicitud / trama} & \textbf{Predicción} & \textbf{Observación} \\ \midrule
\endhead
2 / 6 & 1 & 1 \\
3 / 10 & 2 & 2 \\
4 / 12 & 3 & 3 \\
\bottomrule
\end{longtable}
\endgroup
Alcance: 9 solicitudes en el ejemplo inicial. Una verificación de campo o byte no equivale a una medición independiente; los campos variables conservan sus restricciones y no reciben valores inventados.

\Needspace{9\baselineskip}
\subsubsection{scamper / trace / tcp}
Ejemplo: Valparaíso. Entrenamiento: scamper\_\allowbreak{}20261026T204638\_\allowbreak{}04ca1c5e; contraste: rondas 2 y 3. Las reglas completas de los tres destinos están en el archivo de predicciones enlazado en el anexo.

\begingroup
\fontsize{9}{11}\selectfont
\addtocounter{table}{-1}
\begin{longtable}{p{3.1cm}p{5.8cm}p{5.1cm}}
\toprule
\textbf{Campo} & \textbf{Regla en ronda 1} & \textbf{Contraste posterior} \\ \midrule
\endfirsthead
\toprule
\textbf{Campo} & \textbf{Regla en ronda 1} & \textbf{Contraste posterior} \\ \midrule
\endhead
ip.ttl & incremento 1 cada 2 paquetes & r2: cumple; r3: cumple \\
ip.id & incremento 1 por paquete & r2: cumple; r3: cumple \\
ip.len & constante 40 & r2: cumple; r3: cumple \\
ip.proto & constante 6 & r2: cumple; r3: cumple \\
tcp.srcport & constante 32769 & r2: cumple; r3: cumple \\
tcp.dstport & constante 80 & r2: cumple; r3: cumple \\
tcp.flags & constante 0x0002 & r2: cumple; r3: cumple \\
tcp.seq\_\allowbreak{}raw & constante 0 & r2: cumple; r3: cumple \\
tcp.ack\_\allowbreak{}raw & constante 0 & r2: cumple; r3: cumple \\
tcp.len & constante 0 & r2: cumple; r3: cumple \\
tcp.checksum & constante 0x9922 & r2: cumple; r3: cumple \\
\bottomrule
\end{longtable}
\endgroup
Payload: [0] bytes. No hay datos de aplicación en las sondas TCP observadas.

Predicción concreta en ronda 2, campo ip.ttl: la primera solicitud tiene valor 1. Con la regla aprendida y esa base cuando corresponde, se predicen las solicitudes 2 a 4 antes de contrastarlas con sus valores observados.

\begingroup
\fontsize{9}{11}\selectfont
\addtocounter{table}{-1}
\begin{longtable}{p{3.2cm}p{5.1cm}p{5.1cm}}
\toprule
\textbf{Solicitud / trama} & \textbf{Predicción} & \textbf{Observación} \\ \midrule
\endfirsthead
\toprule
\textbf{Solicitud / trama} & \textbf{Predicción} & \textbf{Observación} \\ \midrule
\endhead
2 / 30 & 1 & 1 \\
3 / 41 & 2 & 2 \\
4 / 54 & 2 & 2 \\
\bottomrule
\end{longtable}
\endgroup
Alcance: 10 solicitudes en el ejemplo inicial. Una verificación de campo o byte no equivale a una medición independiente; los campos variables conservan sus restricciones y no reciben valores inventados.

\Needspace{9\baselineskip}
\subsubsection{scamper / trace / tcp-ack}
Ejemplo: Valparaíso. Entrenamiento: scamper\_\allowbreak{}20261026T204732\_\allowbreak{}c33e1983; contraste: rondas 2 y 3. Las reglas completas de los tres destinos están en el archivo de predicciones enlazado en el anexo.

\begingroup
\fontsize{9}{11}\selectfont
\addtocounter{table}{-1}
\begin{longtable}{p{3.1cm}p{5.8cm}p{5.1cm}}
\toprule
\textbf{Campo} & \textbf{Regla en ronda 1} & \textbf{Contraste posterior} \\ \midrule
\endfirsthead
\toprule
\textbf{Campo} & \textbf{Regla en ronda 1} & \textbf{Contraste posterior} \\ \midrule
\endhead
ip.ttl & incremento 1 cada 2 paquetes & r2: cumple; r3: cumple \\
ip.id & incremento 1 por paquete & r2: cumple; r3: cumple \\
ip.len & constante 40 & r2: cumple; r3: cumple \\
ip.proto & constante 6 & r2: cumple; r3: cumple \\
tcp.srcport & constante 32769 & r2: cumple; r3: cumple \\
tcp.dstport & constante 80 & r2: cumple; r3: cumple \\
tcp.flags & constante 0x0010 & r2: cumple; r3: cumple \\
tcp.seq\_\allowbreak{}raw & constante 0 & r2: cumple; r3: cumple \\
tcp.ack\_\allowbreak{}raw & constante 0 & r2: cumple; r3: cumple \\
tcp.len & constante 0 & r2: cumple; r3: cumple \\
tcp.checksum & constante 0x9914 & r2: cumple; r3: cumple \\
\bottomrule
\end{longtable}
\endgroup
Payload: [0] bytes. No hay datos de aplicación en las sondas TCP observadas.

Predicción concreta en ronda 2, campo ip.ttl: la primera solicitud tiene valor 1. Con la regla aprendida y esa base cuando corresponde, se predicen las solicitudes 2 a 4 antes de contrastarlas con sus valores observados.

\begingroup
\fontsize{9}{11}\selectfont
\addtocounter{table}{-1}
\begin{longtable}{p{3.2cm}p{5.1cm}p{5.1cm}}
\toprule
\textbf{Solicitud / trama} & \textbf{Predicción} & \textbf{Observación} \\ \midrule
\endfirsthead
\toprule
\textbf{Solicitud / trama} & \textbf{Predicción} & \textbf{Observación} \\ \midrule
\endhead
2 / 4 & 1 & 1 \\
3 / 8 & 2 & 2 \\
4 / 10 & 2 & 2 \\
\bottomrule
\end{longtable}
\endgroup
Alcance: 10 solicitudes en el ejemplo inicial. Una verificación de campo o byte no equivale a una medición independiente; los campos variables conservan sus restricciones y no reciben valores inventados.

